\section{Discussion and Limitations}

The four methods form a fairly clear order. Our method reaches 0.265 exact-match accuracy and 0.244 macro-F1, which is close to random four-class performance. This result shows that framing, FFT, local blocks, and normalized correlation can form a complete detection pipeline, but they are not enough for stable four-class discrimination. The intermediate examples explain why: for samples such as \texttt{get} and \texttt{boy}, the four FFT correlation scores are often very close, which means that the maximum-energy block and the one-sided magnitude-spectrum template are not sufficient to separate /g/, /b/, and /d/ reliably.

The AI reimplementation raises exact-match accuracy to 0.503 and macro-F1 to 0.497. This improvement does not come from a black-box classifier. It comes from a more suitable intermediate representation. Log-mel compression smooths spectral shape, onset-centered local windows reduce interference from the vowel in the rest of the word, positive-negative template contrast suppresses energy shared across classes, and autocorrelation adds time-structure information that helps distinguish sustained frication from short stop releases. The class-wise results follow the same explanation. The F1 of /z/ reaches 0.737 because its frication lasts longer and stays more stable. The F1 of /b/ is only 0.359, which shows that a weak release and the following vowel still make it easy to confuse /b/ with /g/ and /d/.

The MLP and CNN are added to examine whether learned models can go beyond the AI reimplementation. The MLP reaches 0.485 exact-match accuracy, which is lower than the AI reimplementation, so whole-word MFCC statistics remain too coarse and average in variation from the word-final vowel and the speaker. The CNN reaches 0.559 exact-match accuracy, which is the best result in the current comparison. It is much stronger on /d/ and /z/, which suggests that after the dataset expands to 300 words and 6796 utterances, the convolutional model does learn some local patterns from the full log-mel spectrogram. Even so, the epoch ablation also shows that the highest development-set macro-F1 at 30 or 40 epochs does not produce the best test accuracy. The best test result still appears at 25 epochs, so the gain of the CNN remains partly dependent on training variance and the data split.

For this reason, ranking the methods by the highest score alone does not tell the whole story. Our method has the lowest score but the clearest interpretation. The AI reimplementation is not the top-scoring method, but it gives the clearest explanation of why a better spectral representation and a better detection structure help. The CNN has the highest score, but its interpretability is the weakest. The comparison also shows that label-wise accuracy is inflated by true negatives under one-hot four-class decoding, so exact-match accuracy, macro-F1, and per-class F1 deserve more attention.

The current study still has several limits. The input remains whole-word audio rather than manually segmented phoneme clips, so our method must rely on sliding-window search. The release cues of /g/, /b/, and /d/ are short, and even small alignment error can bury the template evidence under the following vowel. The current AI reimplementation is still a template method, so it does not use phone-level alignment and does not explicitly model speaker or channel variation. The current CNN is only a light comparison model. It does not include data augmentation, pretrained speech models, or forced alignment, so it should not be treated as an upper bound for deep speech models.
